% ============================================================
%  Übungsaufgaben Template (my_dhbw Stil, uebung_*.tex)
%  Header "Übungsblatt", grün eingefärbte <details>-Lösungen.
%  Renderer: pdflatex (zweimal)
% ============================================================
\documentclass[11pt, a4paper]{article}

\usepackage[ngerman]{babel}
\usepackage{fontspec}
\setmainfont[
  Path=/usr/share/fonts/TTF/,
  Extension=.ttf,
  BoldFont=DejaVuSans-Bold,
  ItalicFont=DejaVuSans-Oblique,
  BoldItalicFont=DejaVuSans-BoldOblique
]{DejaVuSans}
\setsansfont[
  Path=/usr/share/fonts/TTF/,
  Extension=.ttf,
  BoldFont=DejaVuSans-Bold,
  ItalicFont=DejaVuSans-Oblique,
  BoldItalicFont=DejaVuSans-BoldOblique
]{DejaVuSans}
\setmonofont[Path=/usr/share/fonts/TTF/,Extension=.ttf]{DejaVuSansMono}
\usepackage[top=2cm, bottom=2cm, left=2.4cm, right=2.4cm, footskip=1cm]{geometry}
\usepackage{amsmath, amssymb}
\usepackage{xcolor}
\usepackage{enumitem}
\usepackage{fancyhdr}
\usepackage{lastpage}
\usepackage{tabularx}
\usepackage{booktabs}
\usepackage{microtype}
\usepackage{parskip}

\usepackage{array}
\usepackage{longtable}
\usepackage{ltablex}
\keepXColumns
\usepackage{colortbl}
\usepackage[most,breakable]{tcolorbox}
\usepackage{listings}
\usepackage[normalem]{ulem}
\usepackage{pdflscape}
\usepackage{titlesec}
\usepackage[colorlinks=true, linkcolor=accent, urlcolor=accent]{hyperref}

\renewcommand{\familydefault}{\sfdefault}

% ── Theme ─────────────────────────────────────────────────────
\definecolor{accent}{RGB}{0, 60, 113}
\definecolor{primaryblue}{RGB}{0, 60, 113}
\definecolor{loesung}{RGB}{0, 100, 0}
\definecolor{codebg}{RGB}{245, 245, 245}
\definecolor{figurebg}{RGB}{232, 241, 255}
\definecolor{tablehintbg}{RGB}{240, 255, 240}
\definecolor{solutionbg}{RGB}{240, 252, 240}
\definecolor{rulegray}{RGB}{180, 180, 180}
\definecolor{tableheadbg}{RGB}{0, 60, 113}

% ── Header/Footer ─────────────────────────────────────────────
\pagestyle{fancy}
\fancyhf{}
\renewcommand{\headrulewidth}{0.4pt}
\fancyhead[L]{\footnotesize\color{accent}\nichttitle}
\fancyhead[R]{\footnotesize\color{accent}\nichtsubtitle}
\fancyfoot[R]{\footnotesize\color{accent}Blatt \thepage\,/\,\pageref{LastPage}}

% ── Typografie ────────────────────────────────────────────────
\titleformat{\section}
    {\color{accent}\large\bfseries}{}{0pt}{}
    [\vspace{-4pt}\color{accent}\rule{\linewidth}{1.5pt}]
\titleformat{\subsection}
    {\normalsize\bfseries\color{accent!85!black}}{}{0pt}{}{}
\titleformat{\subsubsection}
    {\small\bfseries\color{accent!70!black}}{}{0pt}{}{}
\titlespacing*{\section}{0pt}{10pt}{3pt}
\titlespacing*{\subsection}{0pt}{6pt}{2pt}
\titlespacing*{\subsubsection}{0pt}{4pt}{1pt}

\setlength{\parindent}{0pt}
\setlength{\parskip}{6pt}
\renewcommand{\arraystretch}{1.2}
\setlist[itemize]{noitemsep, topsep=3pt, parsep=0pt, leftmargin=1.4em}
\setlist[enumerate]{noitemsep, topsep=3pt, parsep=0pt, leftmargin=1.6em}

% ── Boxen: solutionbox = Lösungsbox in grün ──────────────────
\newtcolorbox{figurebox}[1]{
  colback=figurebg, colframe=accent!50,
  title={Abbildung: #1}, fonttitle=\small\bfseries,
  breakable, before skip=6pt, after skip=6pt
}
\newtcolorbox{tablehintbox}[1]{
  colback=tablehintbg, colframe=green!50!black!50,
  title={Tabelle: #1}, fonttitle=\small\bfseries,
  breakable, before skip=6pt, after skip=6pt
}
\newtcolorbox{solutionbox}[1]{
  enhanced,
  colback=solutionbg, colframe=loesung,
  title={#1}, fonttitle=\small\bfseries,
  coltitle=white, colbacktitle=loesung,
  breakable, before skip=8pt, after skip=8pt
}

\lstset{
  basicstyle=\ttfamily\small,
  backgroundcolor=\color{codebg},
  breaklines=true,
  frame=single, framerule=0pt,
  xleftmargin=4pt, xrightmargin=4pt,
  aboveskip=6pt, belowskip=6pt,
  keepspaces=true
}

\providecommand{\nichttitle}{Übungsblatt}
\providecommand{\nichtsubtitle}{}

\renewcommand{\nichttitle}{Optimierungsverfahren}
\renewcommand{\nichtsubtitle}{Übungsblatt 1}


\begin{document}

\begin{center}
{\Large\bfseries\color{accent}\nichttitle}
\ifx\nichtsubtitle\empty\else
  \\[2pt]{\normalsize\color{accent!80!black}\nichtsubtitle}
\fi
\\[2pt]
\color{accent}\rule{0.4\linewidth}{0.8pt}\color{black}
\end{center}
\vspace{4pt}

% ── BODY ──────────────────────────────────────────────────────

\section{Übungsblatt 1 — Grundlagen}


\noindent\textcolor{rulegray}{\rule{\linewidth}{0.4pt}}


\paragraph{Aufgabe 1 — Begriffe und Modellaufbau \textit{(8 Punkte)}}\mbox{}\\[2pt]

Im Kapitel werden die Modellbestandteile (Variable, Parameter, Zielfunktion, Nebenbedingung) sowie die Mengen X, X\_z, Y eingeführt.


a) (4 P) Erkläre kurz die Begriffe \textbf{Variable} und \textbf{Parameter} sowie die Mengen \textbf{X} und \textbf{X\_z} und gib für jede der vier Begriffe ein Beispiel aus dem Getränkedose-Modell.


b) (4 P) Was ist die unmittelbare Konsequenz für ein Optimierungsproblem, wenn X\_z = ∅ ist? Begründe und nenne eine konkrete Modellierungsentscheidung im Getränkedose-Modell, die — isoliert verändert — zu X\_z = ∅ führen würde.


\textbf{Lösung:}


a)

\begin{itemize}
  \item \textbf{Variable}: gezielt veränderbare Entscheidungsgröße (Komponenten von x⃗). Getränkedose: Durchmesser d und Höhe h.
  \item \textbf{Parameter}: Umgebungsgröße, die im Modell festliegt und nicht optimiert wird. Getränkedose: Sollvolumen V = 0.5 l und Materialpreis C = 0.01 €/cm².
  \item \textbf{X}: Suchraum, Menge aller im Variablentyp möglichen x⃗. Getränkedose: X = ℝ² (für (d,h)).
  \item \textbf{X\_z}: zulässiger Bereich, X\_z ⊆ X — alle x⃗, die zusätzlich alle Restriktionen erfüllen. Getränkedose: alle (d,h) mit Volumengleichung πd²h/4 = 500, h ≥ 2d, 6 ≤ d ≤ 9 und 12 ≤ h ≤ 18.
\end{itemize}

b) X\_z = ∅ bedeutet: kein einziger Kandidat ist zulässig — das Problem ist \textbf{nicht lösbar} und das Modell muss überarbeitet werden (Restriktion lockern, Variablen ergänzen, Modellfehler suchen).


Beispiel im Getränkedose-Modell: Würde man die Box-Schranke zu 8 ≤ d ≤ 9 verschärfen (statt 6 ≤ d ≤ 9), müsste die Volumengleichung πd²h/4 = 500 erzwingen: bei d = 8 wäre h = 2000/(π·64) ≈ 9.95 cm — verletzt h ≥ 12; bei d = 9 wäre h ≈ 7.86 — auch verletzt. Es existiert kein zulässiges (d,h) mehr, also X\_z = ∅.



\noindent\textcolor{rulegray}{\rule{\linewidth}{0.4pt}}


\paragraph{Aufgabe 2 — Modellierung Solar-Anlage (Transfer) \textit{(22 Punkte)}}\mbox{}\\[2pt]

Eine Hausverwaltung plant eine Solar-Anlage auf einem Flachdach. Pro Modul wird die Jahres-Energieausbeute (in kWh/Jahr) angesetzt zu


e(θ) = 420 − 0.35 · (θ − 30)²,


wobei θ ∈ [0°, 70°] der Aufstellwinkel ist. Pro Modul fallen zudem fixe Jahres-Kosten von k = 180 €/Jahr (Reinigung, Wechselrichter-Anteil, Versicherung) an. Die Anlage darf höchstens 24 Module umfassen, der Strom-Erlös liegt bei p = 0.32 €/kWh. Es soll der Jahresgewinn maximiert werden. Die Modulanzahl n wird kontinuierlich modelliert; θ ist gemeinsamer Aufstellwinkel aller Module. Vertraglich muss die Anlage mindestens 5000 kWh/Jahr liefern.


a) (6 P) Identifiziere \textbf{Variablen} (x⃗), \textbf{Parameter} (p⃗) und benenne \textbf{Zielfunktion}, \textbf{Such- und Zielraum} sowie \textbf{zulässigen Bereich} (in Worten).


b) (10 P) Stelle das Modell vollständig in der \textbf{Standard-Notation} des Kapitels (min f(x), h\_k(x) = 0, g\_j(x) ≤ 0, Box-Constraints) auf. Achte darauf, dass eine Maximierung nicht direkt erlaubt ist und alle Ungleichungen in der Form g\_j ≤ 0 stehen.


c) (6 P) Welcher Punkt (n, θ) wäre dein erster Vermutungs-Kandidat für ein lokales Optimum? Begründe \textbf{ohne} stationäre Punkte explizit zu berechnen — argumentiere allein aus der Struktur von e(θ) und den Restriktionen.


\textbf{Lösung:}


a)

\begin{itemize}
  \item \textbf{Variablen}: x⃗ = (n, θ).
  \item \textbf{Parameter}: p⃗ enthält k = 180, p = 0.32, n\_max = 24, θ\_min = 0, θ\_max = 70, E\_min = 5000 sowie die Strukturkonstanten 420, 0.35 und 30 in e(θ).
  \item \textbf{Zielfunktion}: Jahresgewinn = Erlös − Kosten = n · p · e(θ) − n · k = n · (p · e(θ) − k).
  \item \textbf{Suchraum}: X = ℝ². \textbf{Zielraum}: Y = ℝ (Einzelziel).
  \item \textbf{Zulässiger Bereich}: X\_z = \{ (n, θ) ∈ ℝ² : 0 ≤ n ≤ 24, 0 ≤ θ ≤ 70, n · e(θ) ≥ 5000 \}.
\end{itemize}

b) Maximierung → Minimierung durch Vorzeichenwechsel; ≥ → ≤ durch Multiplikation mit −1.


\begin{lstlisting}
minimiere   f(n, θ) = − n · ( 0.32 · (420 − 0.35·(θ−30)²) − 180 )

mit         g_1(n, θ) = 5000 − n · (420 − 0.35·(θ−30)²) ≤ 0
            (keine Gleichungs-Nebenbedingungen, l = 0)

Box:        0  ≤ n ≤ 24
            0  ≤ θ ≤ 70
\end{lstlisting}


c) e(θ) ist eine nach unten geöffnete Parabel mit Scheitel bei θ = 30°, also e\_max = 420 kWh; θ = 30° liegt im Box-Intervall [0, 70]. Für jedes feste n ist der Gewinn streng monoton in e(θ), also ist θ = 30° unabhängig von n optimal. Bei θ = 30° ist p · e − k = 0.32 · 420 − 180 = −45.6 €/Modul: pro Modul wird Verlust gemacht. Die Zielfunktion fällt deshalb monoton in n — n wäre so klein wie möglich zu wählen. Die Mindestertrags-Restriktion n · 420 ≥ 5000 erzwingt n ≥ 5000/420 ≈ 11.91 und wird also \textbf{aktiv}.


Vermutungs-Kandidat: (n\textit{, θ}) ≈ (11.91, 30°). Das Optimum liegt auf dem \textbf{Rand} des zulässigen Bereichs (Nebenbedingung aktiv) — eine Beobachtung, die für KKT/Lagrange-Methoden später wichtig wird. Der Gewinn ist negativ: das Modell zeigt, dass die Anlage unter den Annahmen unwirtschaftlich ist.



\noindent\textcolor{rulegray}{\rule{\linewidth}{0.4pt}}


\paragraph{Aufgabe 3 — Getränkedose: Zulässigkeit \& Zielfunktion \textit{(18 Punkte)}}\mbox{}\\[2pt]

Gegeben sei das Lernzettel-Modell der Getränkedose:


\begin{lstlisting}
min  f(d, h) = C · π · d · h          (mit C = 0.01 €/cm²)
mit  h_1: π · d² · h / 4 − 500 = 0
     g_1: 2d − h ≤ 0
     6 ≤ d ≤ 9 ;  12 ≤ h ≤ 18
\end{lstlisting}


a) (4 P) Prüfe für jeden der drei Kandidaten \textbf{alle} Restriktionen (h\_1, g\_1, beide Boxen) und gib jeweils „erfüllt" oder „verletzt" an.



{\small
\begin{tabularx}{\linewidth}{XXX}
\hline
\rowcolor{tableheadbg}
{\color{white}\textbf{Kandidat}} & {\color{white}\textbf{d (cm)}} & {\color{white}\textbf{h (cm)}} \\[2pt]
\hline
\endhead
\hline
\endfoot
K1 & 7.0 & 13.0 \\
K2 & 6.5 & 15.07 \\
K3 & 6.0 & 17.68 \\
\end{tabularx}
}


b) (6 P) Eliminiere h aus h\_1 = 0 und drücke f auf der Volumen-Mannigfaltigkeit als Funktion \textbf{nur von d} aus. Berechne f für jeden in (a) als zulässig identifizierten Kandidaten.


c) (8 P) Bestimme den maximal zulässigen Wert d* von d, für den (d, h) auf der Volumen-Mannigfaltigkeit zusätzlich g\_1 erfüllt. Welche Material-Kosten ergeben sich an dieser Stelle? Welcher Kandidat aus (a) liegt diesem Wert am nächsten und ist damit der billigste der drei?


\textbf{Lösung:}


a) Aus h\_1 = 0 folgt d² · h = 2000/π ≈ 636.62.



{\small
\begin{tabularx}{\linewidth}{XXXXXX}
\hline
\rowcolor{tableheadbg}
{\color{white}\textbf{Kandidat}} & {\color{white}\textbf{d²·h}} & {\color{white}\textbf{h\_1 (≈ 0?)}} & {\color{white}\textbf{g\_1: h ≥ 2d?}} & {\color{white}\textbf{d ∈ [6,9]}} & {\color{white}\textbf{h ∈ [12,18]}} \\[2pt]
\hline
\endhead
\hline
\endfoot
K1 & 49 · 13 = 637.0 & erfüllt & \textbf{verletzt} (13 < 14) & erfüllt & erfüllt \\
K2 & 42.25 · 15.07 = 636.7 & erfüllt & erfüllt (15.07 ≥ 13) & erfüllt & erfüllt \\
K3 & 36 · 17.68 = 636.5 & erfüllt & erfüllt (17.68 ≥ 12) & erfüllt & erfüllt \\
\end{tabularx}
}


K1 ist \textbf{nicht zulässig} (g\_1 verletzt). K2 und K3 sind zulässig.


b) Aus h\_1 = 0 folgt h(d) = 2000/(π · d²). Einsetzen:


f(d, h(d)) = C · π · d · 2000/(π · d²) = C · 2000 / d = 20 / d  (€).


\begin{itemize}
  \item K2 (d = 6.5): f = 20 / 6.5 ≈ \textbf{3.077 €}
  \item K3 (d = 6.0): f = 20 / 6.0 ≈ \textbf{3.333 €}
\end{itemize}

c) g\_1 auf der Volumen-Mannigfaltigkeit:


2d − 2000/(π · d²) ≤ 0  ⇔  2π · d³ ≤ 2000  ⇔  d³ ≤ 1000/π ≈ 318.31


⇒  d ≤ d\textit{ := ³√(1000/π) ≈ }\textit{6.828 cm}*.


Da f(d) = 20/d auf [6, 6.828] streng monoton fällt, wird das Minimum in d\textit{ angenommen. h(d}) = 2000/(π · 6.828²) ≈ 13.66 cm liegt in [12, 18] und (d\textit{, h}) liegt in den Boxen.


f(d\textit{) = 20 / 6.828 ≈ }\textit{2.929 €}*.


K2 (d = 6.5) liegt am dichtesten an d\textit{ ≈ 6.83 und ist mit ≈ 3.077 € }\textit{der billigste der drei zulässigen Kandidaten}*.



\noindent\textcolor{rulegray}{\rule{\linewidth}{0.4pt}}


\paragraph{Aufgabe 4 — Umformung in die Standard-Notation \textit{(12 Punkte)}}\mbox{}\\[2pt]

Ein Industrie-Kunde liefert das folgende Optimierungsproblem (x\_1, x\_2, x\_3 ∈ ℝ):


\begin{lstlisting}
maximiere   3·x_1 + 2·x_2 − x_3
mit         x_1 + x_2          ≥ 4
            x_1     − x_3      = 1
            x_2 + x_3          ≤ 5
            −2 ≤ x_1 ≤ 7
             x_2 frei
             x_3 ≥ 0
\end{lstlisting}


Bringe das Problem in die im Kapitel definierte Standard-Notation


\begin{lstlisting}
minimiere f(x), mit h_k(x) = 0, g_j(x) ≤ 0, x_i^l ≤ x_i ≤ x_i^u.
\end{lstlisting}


Notiere alle h\_k, alle g\_j sowie alle Box-Schranken explizit. Begründe in einem Satz pro Schritt, welche Regel des Kapitels du anwendest.


\textbf{Lösung:}


Schritt 1 — \textbf{Maximierung → Minimierung} durch Vorzeichenwechsel der Zielfunktion:


f(x) = −(3·x\_1 + 2·x\_2 − x\_3) = −3·x\_1 − 2·x\_2 + x\_3.


Schritt 2 — \textbf{Ungleichungen in die Form g\_j ≤ 0}:

\begin{itemize}
  \item „x\_1 + x\_2 ≥ 4" mit −1 multiplizieren  →  g\_1(x) = 4 − x\_1 − x\_2 ≤ 0.
  \item „x\_2 + x\_3 ≤ 5" umstellen  →  g\_2(x) = x\_2 + x\_3 − 5 ≤ 0.
\end{itemize}

Schritt 3 — \textbf{Gleichungen in die Form h\_k = 0}:

\begin{itemize}
  \item „x\_1 − x\_3 = 1"  →  h\_1(x) = x\_1 − x\_3 − 1 = 0.
\end{itemize}

Schritt 4 — \textbf{Box-Schranken} explizit:

\begin{itemize}
  \item x\_1: −2 ≤ x\_1 ≤ 7.
  \item x\_2 frei: x\_2\textasciicircum{}l = −∞, x\_2\textasciicircum{}u = +∞.
  \item x\_3 ≥ 0: x\_3\textasciicircum{}l = 0, x\_3\textasciicircum{}u = +∞.
\end{itemize}

\textbf{Standard-Form:}


\begin{lstlisting}
minimiere   f(x) = −3·x_1 − 2·x_2 + x_3
mit         h_1 :  x_1 − x_3 − 1 = 0
            g_1 :  4 − x_1 − x_2 ≤ 0
            g_2 :  x_2 + x_3 − 5 ≤ 0
Box:        −2 ≤ x_1 ≤ 7
            −∞ < x_2 < +∞
             0 ≤ x_3 < +∞
\end{lstlisting}



\noindent\textcolor{rulegray}{\rule{\linewidth}{0.4pt}}


\paragraph{Aufgabe 5 — Modellanalyse: Fehler finden \textit{(15 Punkte)}}\mbox{}\\[2pt]

Ein Mitstudent reicht für die Aufgabe „baue eine quaderförmige Lagerbox mit Breite b, Tiefe t, Höhe h; minimiere Materialbedarf (Oberfläche), Volumen mindestens 1 m³, Tiefe gleich Breite, alle Maße zwischen 0.4 m und 1.5 m" folgende Modellierung ein:


\begin{lstlisting}
maximiere   f(b, t, h, ρ) = 2·(b·t + b·h + t·h)·ρ

mit         h_1 : b · t · h ≥ 1
            h_2 : b − t = 0
            g_1 : b ≤ 1.5
            g_2 : t ≤ 1.5
            g_3 : h ≤ 1.5

            b, t, h ∈ ℝ ;  ρ = 7850 kg/m³
\end{lstlisting}


a) (9 P) Finde \textbf{mindestens vier} Modellfehler bezogen auf die Standard-Notation und die Begriffsdefinitionen aus dem Kapitel. Erkläre jeden Fehler in einem Satz und gib jeweils die korrekte Formulierung an.


b) (6 P) Schreibe das vollständig korrigierte Modell in der Standard-Notation auf.


\textbf{Lösung:}


a) Fehler-Liste (vier reichen — hier sechs benannt):


\begin{enumerate}
  \item \textbf{Maximierung statt Minimierung}: Die Aufgabe verlangt den Materialbedarf zu \textit{minimieren}. Korrektur: \texttt{minimiere f(...)}.
  \item \textbf{Parameter als Variable}: ρ ist eine Materialkonstante (7850 kg/m³) — also ein Parameter. Es darf nicht als Argument in f auftauchen. Korrektur: f(b, t, h) mit ρ als externem Parameter.
  \item \textbf{Zielgröße unscharf}: Wenn das Ziel die \textit{Oberfläche} (Materialbedarf) ist, gehört ρ überhaupt nicht in f, da sie das Optimum nur multiplikativ skaliert. Korrektur: f(b, t, h) = 2 · (bt + bh + th).
  \item \textbf{„h\_1: b·t·h ≥ 1" ist eine Ungleichung, nicht Gleichung}: h\_k bezeichnen im Lernzettel ausschließlich Gleichungen. Außerdem steht ≥ statt ≤. Korrektur: g-Restriktion in Form g\_j ≤ 0:  g\_1 : 1 − b · t · h ≤ 0.
  \item \textbf{Untere Box-Schranken fehlen}: Die Bedingung „mindestens 0.4 m" ist im Modell nicht abgebildet. Korrektur: 0.4 ≤ b, t, h ≤ 1.5.
  \item \textbf{Box-Constraints stehen als g\_j}: Die oberen Schranken b ≤ 1.5 etc. gehören laut Lernzettel in die Box-Section x\_i\textasciicircum{}l ≤ x\_i ≤ x\_i\textasciicircum{}u, nicht in g\_j — sonst wird die Standard-Struktur unsauber.
\end{enumerate}

b) Korrigiertes Modell:


\begin{lstlisting}
minimiere   f(b, t, h) = 2 · (b·t + b·h + t·h)
            (ρ als konstanter Skalierungsfaktor weggelassen, da
             optimumsneutral; ρ bleibt externer Parameter.)

mit         h_1 :  b − t = 0
            g_1 :  1 − b · t · h ≤ 0

Box:        0.4 ≤ b ≤ 1.5
            0.4 ≤ t ≤ 1.5
            0.4 ≤ h ≤ 1.5

Parameter:  ρ = 7850 kg/m³ (nicht Teil der Optimierung)
\end{lstlisting}





% ──────────────────────────────────────────────────────────────

\end{document}
